\documentclass[11pt, a4paper]{article}

\usepackage[margin=2.5cm]{geometry}
\usepackage[utf8]{inputenc}
\usepackage{hyperref}
\usepackage{url}
\usepackage{booktabs}
\usepackage{amsfonts}
\usepackage{xcolor}
\usepackage{amsmath}
\usepackage{amssymb}
\usepackage{amsthm}
\usepackage{graphicx}
\usepackage{subcaption}
\usepackage{algorithm}
\usepackage{algorithmic}

\newtheorem{proposition}{Proposition}

\title{\textbf{Integrating Kolmogorov-Arnold Networks into DreamerV3:\\
A Principled Approach via Binary Latent Representations}}

\author{Przemek Sekula}

\date{}

\begin{document}

\maketitle

\begin{abstract}
World model-based reinforcement learning has achieved strong sample efficiency
through the use of compact latent representations and model-based planning.
DreamerV3~\cite{hafner2023dreamerv3} is a leading example, using a Recurrent
State-Space Model (RSSM) with discrete categorical latent variables. We
investigate replacing the standard multilayer perceptron (MLP) networks in
DreamerV3's actor-critic module with Kolmogorov-Arnold Networks (KANs), which
offer potential benefits for interpretability and, tentatively, for resilience
to catastrophic forgetting. Direct substitution of MLPs with KANs inside
DreamerV3's actor-critic faces two distinct obstacles, however. The first is
practical: KAN spline evaluation scales multiplicatively with input dimension,
the actor-critic receives 1536-dimensional latent features, and it is evaluated
at every step of every imagined rollout, making the computational overhead
severe. The second is a knowledge gap: prior work had not identified a
principled, Dreamer-specific way to alleviate this cost. We address both by
identifying an exact structural property of Dreamer's latent representations.
The stochastic RSSM state is a binary one-hot vector, and for any such binary
input the first KAN layer is mathematically equivalent to a standard affine
transformation---spline expressivity is simply not available at two evaluation
points. This motivates \emph{Sparse-KAN}, a hybrid architecture that replaces
the first KAN layer with a dense linear layer, exploiting the equivalence
exactly, while applying full KAN layers to subsequent continuous representations
where splines can express nontrivial functions. We validate this approach on
the Atari 100k benchmark and demonstrate that the hybrid architecture
successfully learns across a range of Atari games.
\end{abstract}

% ============================================================
\section{Introduction}
% ============================================================

Model-based reinforcement learning (MBRL) has emerged as a powerful paradigm
for achieving sample-efficient learning in complex sequential decision-making
tasks. By learning a world model that predicts how the environment evolves in
response to actions, an agent can generate imagined experience in latent space
and use it to improve its policy without additional interaction with the
environment. DreamerV3~\cite{hafner2023dreamerv3} is one of the most successful
MBRL methods, demonstrating strong performance across domains as diverse as
Atari games, continuous control, and 3-D environments using a single
hyperparameter configuration. Its key architectural components are a Recurrent
State-Space Model (RSSM) with discrete categorical latent variables, a
convolutional encoder and decoder for visual observations, and a model-free
actor-critic that operates entirely in the learned latent space.

Kolmogorov-Arnold Networks (KANs) were recently proposed by Liu et
al.~\cite{liu2024kan} as a principled alternative to multilayer perceptrons
(MLPs). Rather than learning fixed nonlinearities applied to linear projections
of the input, KAN layers learn \emph{univariate activation functions}---realized
as learnable B-spline curves---separately for each connection in the network.
This design has two compelling properties that motivate their use in
reinforcement learning. First, the learned per-edge functions are directly
interpretable as plots of input-output relationships, offering a path toward
explainable neural decision-making. Second, because the spline functions are
local (each B-spline basis function has compact support), updates to one part
of the input space do not necessarily disturb representations elsewhere. This
locality may improve the network's resilience to catastrophic forgetting, a
well-known failure mode in continual and non-stationary learning settings such
as reinforcement learning, though empirical evidence for this advantage is mixed
and appears task-dependent~\cite{kanvmlp2024, kanforgetting2024}.

Despite their promise, KANs face a significant practical obstacle: they are
substantially more expensive to evaluate than equivalent MLPs. The spline
computation introduces a multiplicative overhead proportional to the grid size
and the number of input features, making direct substitution of MLP layers with
KAN layers in high-dimensional networks computationally prohibitive. This
challenge is especially acute in world model-based RL, where the actor and
critic networks receive features of dimension $d_{\text{stoch}} +
d_{\text{deter}} = 1024 + 512 = 1536$ in the default DreamerV3 configuration,
and must be evaluated for every imagined transition over a planning horizon of
15 steps. Prior work integrating KANs into world-model-based RL has addressed
this cost through implementation-level adaptations---vectorization, fixed grids,
simplified initialization---but has not derived a principled justification
rooted in the specific structure of Dreamer's latent representations.

This paper closes both gaps simultaneously. Our key observation is that the
RSSM stochastic state has a specific binary structure that renders the first KAN
layer redundant. Concretely, the RSSM represents the stochastic
component of the latent state as a collection of independent categorical
variables: 32 groups, each represented as a one-hot vector of dimension 32,
yielding a 1024-dimensional binary vector. The actor and critic networks
receive this flattened one-hot representation as a dominant part of their
input. We show analytically that, for inputs restricted to $\{0, 1\}^n$, a
KAN linear layer is exactly equivalent to a standard affine transformation: the
learnable spline functions collapse to linear functions of the binary input, and
the net computation reduces to a weight matrix multiplication plus a bias term.
This equivalence means that using a full KAN layer at the input of the
actor-critic network wastes capacity---the expressive power of splines is not
utilized, yet the spline computation is still performed.

Motivated by this observation, we design \emph{Sparse-KAN}, a hybrid
architecture in which the first layer of the actor and critic is replaced by a
standard dense linear layer (exploiting the mathematical equivalence), and all
subsequent layers are KAN layers (where inputs are no longer binary, so splines
can express nontrivial functions). This design recovers the full expressive
power of KANs in the intermediate representations without paying for their
computational overhead at the binary input stage.

We evaluate Sparse-KAN on the Atari 100k benchmark~\cite{kaiser2020atari100k},
a standard testbed for sample-efficient visual reinforcement learning in which
agents are limited to $4 \times 10^5$ environment steps. Our experiments show
that the hybrid architecture successfully learns across a diverse set of Atari
games, matching or approaching the performance of the standard MLP-based
DreamerV3 baseline. To the best of our knowledge, this is the first work in Dreamer-style
world model reinforcement learning to exploit the binary one-hot structure of
the RSSM stochastic state and show that the first KAN layer collapses to an
affine map, motivating a hybrid linear-plus-KAN actor-critic.


\section{Related Work}

\paragraph{World model-based reinforcement learning.}
World model-based reinforcement learning has developed into a strong paradigm
for sample-efficient control by learning latent dynamics and improving policies
from imagined rollouts rather than only from direct environment interaction.
Early work on \emph{World Models}~\cite{ha2018world} showed that compact learned simulators can
support downstream policy learning, while PlaNet~\cite{hafner2019planet} introduced recurrent
state-space models with stochastic and deterministic latent components for
planning in latent space. Dreamer~\cite{hafner2019dreamer} then demonstrated that actor and critic
networks can be trained directly on imagined trajectories, avoiding explicit
planning at decision time. This line culminated in DreamerV2~\cite{hafner2021dreamerv2} and
DreamerV3~\cite{hafner2023dreamerv3}, which established discrete latent world models as a highly
effective basis for visual control and broad-domain reinforcement learning. In
particular, DreamerV2 introduced discrete world models for Atari, and DreamerV3
showed that a single world-model agent can perform strongly across diverse
domains with fixed hyperparameters. 

\paragraph{Kolmogorov-Arnold Networks and their efficient variants.}
Kolmogorov-Arnold Networks (KANs) were proposed as an alternative to standard
MLPs in which each input-output connection is parameterized by a learnable
univariate function, typically implemented with spline-based bases, rather than
by a single scalar weight. The original KAN work motivated this design through
interpretability and local functional plasticity, and it triggered a rapidly
growing literature on practical variants. Since the original formulation is
computationally expensive, several follow-up works have focused on efficiency,
including FastKAN~\cite{li2024fastkan}, which approximates spline bases with radial basis
functions, ReLU-KAN~\cite{qiu2024relukan}, which uses matrix-friendly ReLU constructions,
UKAN~\cite{ukan2024}, which addresses bounded-grid limitations, and Efficient-KAN~\cite{blealtan2024},
which provides a more memory-efficient implementation. At the same time, the
empirical evidence on KANs remains mixed: recent careful comparisons suggest
that KANs do not consistently outperform well-matched MLPs outside selected
settings~\cite{kanvmlp2024}, and their advantages for catastrophic forgetting
appear to be task-dependent rather than universal~\cite{kanforgetting2024}. 

\paragraph{KANs in reinforcement learning.}
A growing body of work has already explored KANs as function approximators in
reinforcement learning. Prior studies have evaluated KAN-based PPO in online
control~\cite{kanppo2024} and KAN-based actor-critic architectures in offline reinforcement
learning~\cite{kanoffline2024}, often finding that KANs can achieve performance comparable to MLP
baselines while sometimes using fewer parameters. Other reinforcement-learning
applications have emphasized interpretability, for example in network control
and value-based learning settings. These studies show that the idea of using
KANs in reinforcement learning is no longer novel by itself. Rather, the open
question is whether KANs can be integrated effectively into
\emph{world model-based} reinforcement learning, where actor and critic
networks are evaluated repeatedly during imagination rollouts and computational
overhead becomes especially important. 

\paragraph{Closest prior work.}
The closest prior work we found is \emph{KAN-Dreamer}~\cite{kandreamer2025}, which explicitly studies
the integration of KAN and FastKAN components into DreamerV3. That work
considers replacing MLP-based modules in visual perception, latent prediction,
and behavior learning, and it reports that KAN-based actor-critic variants can
train successfully in DreamerV3. Consequently, a broad claim that the present
paper is the first successful integration of KANs into a world model-based RL
agent would be too strong. However, the contribution of the present paper is
different in emphasis and mechanism. Whereas prior Dreamer/KAN work motivates
feasibility primarily through implementation-oriented choices such as
vectorization, fixed grids, and simplified initialization, our approach is
based on a structural property of Dreamer-style latent states: the stochastic
RSSM component is represented as binary one-hot variables. This observation
leads to a principled result showing that a KAN input layer acting on binary
inputs reduces to an affine map, which in turn motivates replacing the first
KAN layer with a standard dense transformation and reserving KAN layers for
subsequent continuous representations. To the best of our knowledge, this
specific binary-input equivalence and the resulting hybrid linear-plus-KAN
design have not been previously identified in the Dreamer/KAN literature. 

\paragraph{Positioning of the present work.}
Taken together, the literature suggests that the main novelty of the present
paper is not the general use of KANs in reinforcement learning, nor the mere
replacement of MLP components inside DreamerV3. Instead, the contribution lies
in identifying a Dreamer-specific architectural simplification that follows
from the binary one-hot structure of the discrete RSSM latent state, and in
using that result to design a computationally feasible hybrid actor-critic for
world model-based RL. This positions our work at the intersection of efficient
world-model reinforcement learning and hybrid KAN architectures. 

% ============================================================
\section{Methods}
% ============================================================

\subsection{DreamerV3 Background}

DreamerV3 learns a world model of the form:
\begin{align}
  \text{Encoder:}       &\quad z_t \sim q_\phi(z_t \mid o_t, h_t) \\
  \text{Recurrence:}    &\quad h_t = f_\phi(h_{t-1},\, z_{t-1},\, a_{t-1}) \\
  \text{Dynamics:}      &\quad \hat{z}_t \sim p_\phi(\hat{z}_t \mid h_t) \\
  \text{Decoder:}       &\quad \hat{o}_t \sim p_\phi(\hat{o}_t \mid h_t,\, z_t) \\
  \text{Reward / cont.:}&\quad \hat{r}_t,\, \hat{c}_t \sim p_\phi(\cdot \mid h_t,\, z_t)
\end{align}
where $o_t$ is the observation, $a_t$ is the action, $z_t$ is the stochastic
latent state, and $h_t$ is the deterministic recurrent state. The recurrence
$f_\phi$ is a Gated Recurrent Unit (GRU), making $h_t$ the hidden state of the
memory network.

\paragraph{Discrete latent variables.}
The stochastic state $z_t$ is parameterized as a product of $N_s = 32$
independent categorical distributions, each over $N_d = 32$ categories. At
each step, a single category is selected per group, yielding a one-hot
encoding. The full stochastic state is the concatenation of these one-hot
vectors, flattened to $z_t \in \{0,1\}^{1024}$. Training uses the
straight-through estimator to allow gradients to flow through the discrete
sampling operation.

\paragraph{Feature representation.}
The complete state feature used by the actor and critic is the concatenation
\begin{equation}
  \mathbf{f}_t = [\,z_t;\, h_t\,] \in \mathbb{R}^{1536},
\end{equation}
where $z_t \in \{0,1\}^{1024}$ is the binary stochastic state and $h_t \in
\mathbb{R}^{512}$ is the continuous GRU hidden state.

\paragraph{Behavior learning via imagination.}
DreamerV3 trains an actor $\pi_\theta$ and critic $v_\psi$ entirely within the
world model. Starting from an encoded state, the agent imagines a trajectory of
$H = 15$ steps by repeatedly sampling actions from $\pi_\theta$, advancing the
dynamics with $f_\phi$, and evaluating $v_\psi$ to compute bootstrapped
$\lambda$-returns. The actor is trained to maximize these imagined returns; the
critic is trained to predict them. Both networks receive the feature
$\mathbf{f}_t$ as input and are implemented as MLPs in the original DreamerV3.

\subsection{Kolmogorov-Arnold Networks}

\paragraph{Formulation.}
The Kolmogorov-Arnold representation theorem states that any multivariate
continuous function can be expressed as a finite composition of univariate
functions~\cite{kolmogorov1957superpositions}. KANs operationalize this idea by
replacing the learned scalar weights of a standard linear layer with learned
univariate activation functions, one per input-output connection. A single KAN
linear layer with $n$ inputs and $m$ outputs computes:
\begin{equation}
  y_k = \sum_{i=1}^{n} \phi_{ki}(x_i), \quad k = 1, \ldots, m,
\end{equation}
where each $\phi_{ki}: \mathbb{R} \to \mathbb{R}$ is a learnable function. In
practice, each $\phi_{ki}$ is parameterized as the sum of a base activation and
a learnable B-spline~\cite{liu2024kan}:
\begin{equation}
  \phi_{ki}(x) = w^{\text{base}}_{ki} \cdot b(x) + \sum_{j=1}^{G+p} c_{kij} \cdot B_j(x),
  \label{eq:kanlinear}
\end{equation}
where $b(\cdot) = \text{SiLU}(\cdot)$ is the base activation, $\{B_j\}_{j=1}^{G+p}$
is a B-spline basis of order $p$ defined over a uniform grid of $G{+}1$ knots
on the interval $[-1, 1]$, and $\{c_{kij}\}$ are the learnable spline
coefficients. The parameter $w^{\text{base}}_{ki}$ controls the strength of
the residual linear component.

\paragraph{Computational overhead.}
Evaluating \eqref{eq:kanlinear} for a single layer requires computing $n \times
m$ B-spline evaluations, each involving $G + p$ basis functions. For a
high-dimensional input ($n = 1536$ in our setting) and typical grid size
($G = 8$) and spline order ($p = 3$), this introduces an order-of-magnitude
increase in floating-point operations compared to an MLP layer of the same
width. This overhead motivates the architectural optimization described next.

\paragraph{Efficient sparse KAN.}
We build on the efficient KAN implementation of~\cite{blealtan2024}, which
introduces a sparse B-spline computation that skips evaluations for zero entries
in the input tensor, exploiting the assumption $f(0) = 0$ for the spline
component. This sparse formulation is particularly effective for binary inputs,
where the majority of entries are zero.

\subsection{Binary Input Equivalence}

The binary nature of the discrete latent state $z_t$ has an important
consequence for any KAN layer that receives $\mathbf{f}_t$ as input.

\begin{proposition}[KAN linear layer with binary inputs is affine]
  \label{prop:binary-equivalence}
  Let $\mathbf{x} \in \{0,1\}^n$ be a binary input vector. Then the output of
  a KANLinear layer as defined in \eqref{eq:kanlinear} is an affine function of
  $\mathbf{x}$; that is, there exist $W^{\mathrm{eff}} \in \mathbb{R}^{m \times n}$
  and $\mathbf{b}^{\mathrm{eff}} \in \mathbb{R}^m$ such that
  \begin{equation}
    \mathrm{KANLinear}(\mathbf{x}) = W^{\mathrm{eff}}\, \mathbf{x} + \mathbf{b}^{\mathrm{eff}}.
  \end{equation}
\end{proposition}

\begin{proof}
Fix any output unit $k$ and input feature $i$. Since $x_i \in \{0, 1\}$ and
the spline grid is fixed, each B-spline basis function $B_j$ evaluates to one
of two constants at the input: $B_j(0)$ or $B_j(1)$. Likewise, $b(0) =
\mathrm{SiLU}(0) = 0$ and $b(1) = \mathrm{SiLU}(1)$ are constants independent
of the learnable parameters. Using the identity $B_j(x_i) = x_i B_j(1) +
(1 - x_i) B_j(0)$ and $b(x_i) = x_i b(1)$, we expand \eqref{eq:kanlinear}:
\begin{align}
  \phi_{ki}(x_i)
  &= w^{\mathrm{base}}_{ki}\, x_i\, b(1)
   + \sum_{j} c_{kij} \bigl[ x_i B_j(1) + (1 - x_i) B_j(0) \bigr] \notag \\
  &= x_i \underbrace{\Bigl[ w^{\mathrm{base}}_{ki}\, b(1)
       + \textstyle\sum_{j} c_{kij}\bigl(B_j(1) - B_j(0)\bigr) \Bigr]}_{W^{\mathrm{eff}}_{ki}}
   + \underbrace{\sum_{j} c_{kij}\, B_j(0)}_{\beta_{ki}}.
\end{align}
Summing over all input features:
\begin{equation}
  y_k = \sum_{i=1}^{n} \phi_{ki}(x_i)
      = \sum_{i=1}^{n} W^{\mathrm{eff}}_{ki}\, x_i
        + \underbrace{\sum_{i=1}^{n} \beta_{ki}}_{b^{\mathrm{eff}}_k},
\end{equation}
which is the $k$-th component of $W^{\mathrm{eff}}\mathbf{x} + \mathbf{b}^{\mathrm{eff}}$.
\end{proof}

The proposition shows that applying a KAN layer to binary inputs is no more
expressive than applying a standard linear layer: with only two evaluation
points per function ($\{0, 1\}$), the learnable splines cannot represent any
nonlinearity. Any function capacity stored in the spline coefficients is
therefore redundant with the capacity of a linear weight matrix. Consequently,
using a full KAN as the first layer of the actor or critic---when its input
contains a dominant binary component from the discrete latent state---wastes
parameters and computation without providing any representational benefit.

\subsection{Sparse-KAN: Hybrid Architecture}

\paragraph{Architecture.}
Based on Proposition~\ref{prop:binary-equivalence}, we design the Sparse-KAN
module used for both the actor and the critic. Given the input feature
$\mathbf{f}_t \in \mathbb{R}^{1536}$, the network computes:
\begin{equation}
  \mathbf{h}_0 = \mathrm{SiLU}(W_0\, \mathbf{f}_t + \mathbf{b}_0),
  \qquad
  \mathbf{h}_\ell = \mathrm{KAN}_\ell(\mathbf{h}_{\ell-1}), \quad \ell = 1, \ldots, L,
\end{equation}
where $W_0 \in \mathbb{R}^{d_0 \times 1536}$ is a standard trainable weight
matrix and $\mathrm{KAN}_\ell$ are KAN linear layers as defined in
\eqref{eq:kanlinear}. The first transformation is an ordinary dense linear
projection followed by a SiLU nonlinearity; subsequent layers are full KAN
layers that exploit the expressive power of learnable splines over continuous
intermediate representations.

This design is mathematically justified by Proposition~\ref{prop:binary-equivalence}:
replacing the first KAN layer with a dense linear layer is exact for the binary
portion of the input. The continuous portion $h_t$ enters the network jointly
through the same first projection, and nonlinear processing by the spline
functions begins in the second layer once the activations are no longer binary.

\paragraph{Layer widths and spline parameters.}
For Atari 100k experiments we use hidden widths $[d_0, d_1, d_2] = [128, 64,
32]$, B-spline grid size $G = 8$, and spline order $p = 3$ (cubic splines).
The same architecture is used for both actor and critic. The output heads
(action logits for the actor, value distribution logits for the critic) are
standard linear layers applied to $\mathbf{h}_L$.

\paragraph{Action distribution.}
In Atari, actions are discrete; the actor samples from a one-hot categorical
distribution using the REINFORCE gradient estimator~\cite{williams1992simple}.
This means the action vectors that enter the RSSM as previous actions $a_{t-1}$
are also binary one-hot vectors, further motivating the linear first-layer
treatment for inputs to the dynamics model.

\subsection{Training Setup}

\paragraph{Environment and benchmark.}
We evaluate on the Atari 100k benchmark~\cite{kaiser2020atari100k}, in which
agents are trained for $4 \times 10^5$ environment steps ($10^5$ steps with
action repeat 4) and evaluated over 100 episodes. We train on 26 Atari games
using standard preprocessing: grayscale observations, no-op initialization with
up to 30 no-ops, and episode termination on life loss disabled.

\paragraph{World model.}
The world model follows the DreamerV3 default for Atari 100k: a convolutional
encoder and decoder with 32 channels, RSSM with $N_s = 32$ stochastic
categories of dimension $N_d = 32$ (so $z_t \in \{0,1\}^{1024}$) and
deterministic hidden dimension $d_h = 512$, trained with a replay ratio of
1024 updates per environment step. The world model components (encoder, RSSM,
decoder, reward head, continuation head) use standard MLP layers as in the
original DreamerV3, unchanged.

\paragraph{Actor and critic.}
Only the actor and critic networks are replaced with the Sparse-KAN
architecture described above. Isolating the replacement to the behavior module
allows us to measure the effect of KAN on policy learning independently of any
change in world model quality. The actor uses REINFORCE; the critic uses a
symlog-encoded two-hot distribution over returns.

\paragraph{Optimization.}
All parameters are optimized with Adam. The world model uses learning rate
$10^{-4}$; the actor and critic use $3 \times 10^{-5}$. Gradient norms are
clipped at $10^3$ for the world model and $10^2$ for the behavior module. We
use a batch size of 16 sequences of length 64.

\paragraph{Baseline.}
The baseline is the standard DreamerV3 with MLP actor and critic, identical in
all hyperparameters except for the network architecture. Performance is reported
as mean and standard deviation of undiscounted episode returns across seeds and
evaluation episodes.

% ============================================================
\section{Results}
% ============================================================

\subsection{Per-game Performance}

Table~\ref{tab:per_game} reports the human-normalized score (HNS) for each of the 26 Atari~100k
games at the end of training for both the DreamerV3 baseline (MLP actor-critic) and the
proposed Sparse-KAN actor-critic. All other components of the agent---the world model,
encoder, decoder, and training hyperparameters---are identical between the two conditions.

\begin{table}[ht]
\centering
\caption{Human-normalized scores (HNS) at the final checkpoint across 26 Atari~100k games.
  Scores above~1 exceed the reference human level.
  Bold indicates the higher score for each game.
  $^\dagger$Breakout data is unavailable for Sparse-KAN; mean is computed over the remaining 25 games.}
\label{tab:per_game}
\setlength{\tabcolsep}{5pt}
\begin{tabular}{lcc|lcc}
\toprule
Game & \shortstack{DreamerV3\\(baseline)} & \shortstack{Sparse-KAN\\(ours)}
   & Game & \shortstack{DreamerV3\\(baseline)} & \shortstack{Sparse-KAN\\(ours)} \\
\midrule
Alien          & \textbf{0.150} & 0.080 & Kangaroo       & \textbf{1.355} & 0.586 \\
Amidar         & \textbf{0.086} & 0.037 & Krull          & 6.177 & \textbf{7.537} \\
Assault        & 1.060 & \textbf{1.175} & Kung Fu Master & \textbf{1.138} & 1.022 \\
Asterix        & \textbf{0.111} & 0.099 & Ms.\ Pac-Man   & \textbf{0.279} & 0.057 \\
Bank Heist     & \textbf{1.231} & 0.434 & Pong           & 1.093 & \textbf{1.153} \\
Battle Zone    & 0.215 & \textbf{0.407} & Private Eye    & \textbf{$-$0.003} & $-$0.000 \\
Boxing         & 6.905 & \textbf{7.352} & Qbert          & \textbf{0.182} & 0.035 \\
Breakout       & \textbf{0.938} & N/A$^\dagger$ & Road Runner    & 1.870 & \textbf{1.940} \\
Chopper Cmd.   & \textbf{0.199} & 0.040 & Seaquest       & \textbf{0.014} & 0.003 \\
Crazy Climber  & \textbf{2.790} & 2.670 & Up N Down      & \textbf{0.942} & 0.491 \\
Demon Attack   & 0.048 & \textbf{0.062} & & & \\
Freeway        & \textbf{0.216} & 0.000 & & & \\
Frostbite      & \textbf{0.561} & 0.045 & & & \\
Gopher         & 1.118 & \textbf{1.298} & & & \\
Hero           & \textbf{0.169} & 0.079 & & & \\
Jamesbond      & \textbf{1.542} & 1.507 & & & \\
\midrule
\multicolumn{4}{l}{Mean HNS (25 games$^\dagger$)}   & \textbf{1.169} & 1.124 \\
\multicolumn{4}{l}{Median HNS} & \textbf{0.750} & 0.407 \\
\multicolumn{4}{l}{Games $>$ human} & \textbf{11} & 9 \\
\bottomrule
\end{tabular}
\end{table}

\subsection{Comparison with Published Methods}

Table~\ref{tab:aggregate} places our two conditions alongside representative published
methods on Atari~100k, drawn from~\cite{schwarzer2023bbf}.  Because computing the
interquartile mean (IQM) robustly requires per-seed, per-game
score distributions that we do not have for all baselines, we report mean and
median HNS for all methods, using mean as the primary metric.  Published values for
the literature baselines are taken directly from Table~A.1 of~\cite{schwarzer2023bbf}.

\begin{table}[ht]
\centering
\caption{Aggregate human-normalized scores on Atari~100k (26 games), sorted by mean HNS.
  Values for literature methods are taken from~\cite{schwarzer2023bbf};
  values for our two conditions are computed from the final training checkpoint
  averaged across seeds.  Higher is better.
  $^\dagger$Sparse-KAN mean is over 25 games (Breakout data unavailable).}
\label{tab:aggregate}
\begin{tabular}{lcc}
\toprule
Method & Mean HNS & Median HNS \\
\midrule
Random                     & 0.000 & 0.000 \\
DER                        & 0.350 & 0.189 \\
DrQ($\epsilon$)            & 0.465 & 0.313 \\
SPR                        & 0.616 & 0.396 \\
IRIS                       & 1.046 & 0.289 \\
\textbf{Sparse-KAN (ours)}$^\dagger$ & \textbf{1.124} & \textbf{0.407} \\
DreamerV3 (baseline, ours) & 1.169 & 0.750 \\
SR-SPR                     & 1.272 & 0.685 \\
EfficientZero              & 1.945 & 1.116 \\
BBF                        & 2.247 & 0.917 \\
\midrule
Human                      & 1.000 & 1.000 \\
\bottomrule
\end{tabular}
\end{table}

\subsection{Discussion}

\paragraph{The method clearly works.}
Sparse-KAN successfully learns non-trivial policies across all 26 Atari~100k games,
achieving super-human performance on 9 of them.  A mean HNS of~1.12 places it
just above IRIS and just below the DreamerV3 MLP baseline (Table~\ref{tab:aggregate}),
despite using a KAN architecture that was never tuned for performance.  Some of this mean is driven by games where
both conditions dramatically exceed human level (most notably Boxing and Krull, where
spline-based representations appear particularly effective), but the breadth of
learning---positive returns on every game---confirms that the KAN-based actor-critic
is a viable function approximator in the DreamerV3 imagination rollout loop.

\paragraph{Comparison with the MLP baseline.}
Sparse-KAN underperforms the DreamerV3 MLP baseline (mean~1.12 vs.~1.17), a gap
that is visible across most individual games in Table~\ref{tab:per_game}.  This is
not surprising.  The MLP baseline uses the same width configuration that was tuned
for DreamerV3 over many experiments by the original authors, while our KAN
architecture---hidden widths $[128, 64, 32]$ with $G=8$ spline grid and cubic
splines---was selected primarily to be computationally feasible rather than
empirically optimal.  Importantly, the Sparse-KAN architecture still matches or
exceeds the baseline on several games (Kung Fu Master, Up N Down, Jamesbond), which
shows that the KAN layers are capable of learning useful value and policy functions
when the representation is continuous, exactly as the theory in
Section~3.3 predicts.

\paragraph{Computational constraints and room for optimization.}
Due to hardware limitations, we were unable to conduct a systematic search over KAN
architecture hyperparameters.  The hidden-layer widths, spline grid size, spline order,
and number of KAN layers were all fixed without tuning.  Prior comparative studies
suggest that KAN performance is sensitive to these choices~\cite{kanvmlp2024}, and
it is plausible that the current configuration is suboptimal.  In particular, wider
hidden layers or a finer spline grid (larger $G$) may allow the KAN layers to capture
more expressive value and policy functions in the continuous intermediate
representations where splines are beneficial.  We therefore view the 1.12 mean HNS
as a lower bound on what the Sparse-KAN architecture can achieve with proper
hyperparameter tuning, and we expect that the gap to the MLP baseline is reducible---
possibly substantially---with additional compute for architecture search.

% ============================================================
\bibliographystyle{plain}
\bibliography{dreamer}

\end{document}
